\documentclass[11pt]{article}

% Common class-neutral preamble for standalone notes under manuscript/notes/.
% Note entry points all live exactly two directories below manuscript/.

\usepackage[margin=1in]{geometry}
\usepackage{amsmath,amssymb,amsthm,mathtools,bm}
\usepackage{algorithm}
\usepackage{algorithmic}
\usepackage{booktabs,tabularx,multirow,array,longtable}
\usepackage{enumitem}
\usepackage{microtype}
\usepackage[numbers,sort&compress]{natbib}
\usepackage{xcolor}
\usepackage[colorlinks=true,citecolor=blue,linkcolor=black,urlcolor=blue]{hyperref}
\usepackage[capitalize,nameinlink,noabbrev]{cleveref}

\newtheorem{theorem}{Theorem}[section]
\newtheorem{lemma}[theorem]{Lemma}
\newtheorem{proposition}[theorem]{Proposition}
\newtheorem{corollary}[theorem]{Corollary}
\newtheorem{conjecture}[theorem]{Conjecture}
\theoremstyle{definition}
\newtheorem{definition}[theorem]{Definition}
\newtheorem{assumption}[theorem]{Assumption}
\newtheorem{problem}[theorem]{Problem}
\newtheorem{observation}[theorem]{Observation}
\theoremstyle{remark}
\newtheorem{remark}[theorem]{Remark}
\theoremstyle{plain}

% Canonical mathematical notation for the active manuscript.
%
% This file preserves the recurring notation style used in the author's
% NeurIPS 2024 and 2025 papers while excluding unrelated tensor and
% machine-learning boilerplate from those archived command collections.
%
% Prefix convention:
%   \v...  bold vectors
%   \m...  bold matrices
%   \g...  calligraphic sets, graphs, and families
%   \s...  blackboard-bold spaces and number systems

\newcommand{\method}{\textsc{HybridLocalSolver}}

% Font and scalar shorthands retained from the archived manuscripts.
\newcommand{\mc}{\mathcal}
\newcommand{\R}{\mathbb{R}}
\newcommand{\E}{\mathbb{E}}
\newcommand{\G}{\mathbb{G}}
\newcommand{\N}{\mathcal{N}}
% Accuracy namespaces are semantic: do not add a bare \eps alias.
% Degree-normalized residual tolerance of classical APPR is kept distinct from
% the objective-gap and proximal fixed-point tolerances of the RPPR sections.
\newcommand{\epsappr}{\varepsilon_{\mathrm{appr}}}
\newcommand{\epsobj}{\varepsilon_{\mathrm{obj}}}
\newcommand{\epspg}{\varepsilon_{\mathrm{pg}}}
\newcommand{\epsppr}{\varepsilon_{\mathrm{ppr}}}
\newcommand{\epsin}{\varepsilon_{\mathrm{in}}}
\newcommand{\epskkt}{\varepsilon_{\mathrm{kkt}}}
\newcommand{\epsburn}{\varepsilon_{\mathrm{burn}}}
\newcommand{\epssol}{\varepsilon_{\mathrm{sol}}}
\newcommand{\epspprinst}[1]{\varepsilon_{\mathrm{ppr},#1}}
\newcommand{\trans}{\mathsf{T}}
\newcommand{\grad}{\nabla}

% Delimiters and generic helpers.
% Norms and inner products use fixed-size delimiters by default.
% Use an optional size (e.g. \norm[\big]{...}) or * for automatic sizing.
\DeclarePairedDelimiter{\norm}{\lVert}{\rVert}
\newcommand{\abs}[1]{\left\lvert #1 \right\rvert}
\DeclarePairedDelimiterX{\inner}[2]{\langle}{\rangle}{#1, #2}
\newcommand{\ip}[2]{\inner{#1}{#2}}
\newcommand{\card}[1]{\left\lvert #1 \right\rvert}
\newcommand{\set}[1]{\left\{#1\right\}}
\newcommand{\ceil}[1]{\left\lceil #1 \right\rceil}
\newcommand{\floor}[1]{\left\lfloor #1 \right\rfloor}
\newcommand{\comp}[1]{\overline{#1}}
\newcommand{\conj}[1]{\overline{#1}}
\newcommand{\mean}[1]{\overline{#1}}
\newcommand{\bigO}{\mathcal{O}}
\newcommand{\softO}{\widetilde{\mathcal{O}}}
\newcommand{\softOmega}{\widetilde{\Omega}}
\newcommand{\pos}[1]{\left[#1\right]_{+}}

% Bold lowercase vectors.
\newcommand{\va}{\bm{a}}
\newcommand{\vb}{\bm{b}}
\newcommand{\vc}{\bm{c}}
\newcommand{\vd}{\bm{d}}
\newcommand{\ve}{\bm{e}}
\newcommand{\vf}{\bm{f}}
\newcommand{\vg}{\bm{g}}
\newcommand{\vh}{\bm{h}}
\newcommand{\vi}{\bm{i}}
\newcommand{\vj}{\bm{j}}
\newcommand{\vk}{\bm{k}}
\newcommand{\vl}{\bm{l}}
\newcommand{\vm}{\bm{m}}
\newcommand{\vn}{\bm{n}}
\newcommand{\vo}{\bm{o}}
\newcommand{\vp}{\bm{p}}
\newcommand{\vq}{\bm{q}}
\newcommand{\vr}{\bm{r}}
\newcommand{\vs}{\bm{s}}
\newcommand{\vt}{\bm{t}}
\newcommand{\vu}{\bm{u}}
\newcommand{\vv}{\bm{v}}
\newcommand{\vw}{\bm{w}}
\newcommand{\vx}{\bm{x}}
\newcommand{\vy}{\bm{y}}
\newcommand{\vz}{\bm{z}}

% Bold Greek vectors and special vectors.
\newcommand{\valpha}{\bm{\alpha}}
\newcommand{\vbeta}{\bm{\beta}}
\newcommand{\vdelta}{\bm{\delta}}
\newcommand{\vepsilon}{\bm{\epsilon}}
\newcommand{\veta}{\bm{\eta}}
\newcommand{\vgamma}{\bm{\gamma}}
\newcommand{\vlambda}{\bm{\lambda}}
\newcommand{\vmu}{\bm{\mu}}
\newcommand{\vomega}{\bm{\omega}}
\newcommand{\vphi}{\bm{\phi}}
\newcommand{\vpi}{\bm{\pi}}
\newcommand{\vpsi}{\bm{\psi}}
\newcommand{\vrho}{\bm{\rho}}
\newcommand{\vsigma}{\bm{\sigma}}
\newcommand{\vtheta}{\bm{\theta}}
\newcommand{\vxi}{\bm{\xi}}
\newcommand{\vell}{\bm{\ell}}
\newcommand{\vDelta}{\bm{\Delta}}
\newcommand{\vGamma}{\bm{\Gamma}}
\newcommand{\vLambda}{\bm{\Lambda}}
\newcommand{\vPhi}{\bm{\Phi}}
\newcommand{\vSigma}{\bm{\Sigma}}
\newcommand{\vzero}{\bm{0}}
\newcommand{\vone}{\bm{1}}
\newcommand{\one}{\mathbf{1}}
\newcommand{\zero}{\vzero}
\newcommand{\eunit}[1]{\bm{e}_{#1}}
\newcommand{\vDeltax}{\bm{\Delta}\bm{x}}

% Bold uppercase matrices.
\newcommand{\mA}{\bm{A}}
\newcommand{\mB}{\bm{B}}
\newcommand{\mC}{\bm{C}}
\newcommand{\mD}{\bm{D}}
\newcommand{\mE}{\bm{E}}
\newcommand{\mF}{\bm{F}}
\newcommand{\mG}{\bm{G}}
\newcommand{\mH}{\bm{H}}
\newcommand{\mI}{\bm{I}}
\newcommand{\mJ}{\bm{J}}
\newcommand{\mK}{\bm{K}}
\newcommand{\mL}{\bm{L}}
\newcommand{\mM}{\bm{M}}
\newcommand{\mN}{\bm{N}}
\newcommand{\mO}{\bm{O}}
\newcommand{\mP}{\bm{P}}
\newcommand{\mQ}{\bm{Q}}
\newcommand{\mR}{\bm{R}}
\newcommand{\mS}{\bm{S}}
\newcommand{\mT}{\bm{T}}
\newcommand{\mU}{\bm{U}}
\newcommand{\mV}{\bm{V}}
\newcommand{\mW}{\bm{W}}
\newcommand{\mX}{\bm{X}}
\newcommand{\mY}{\bm{Y}}
\newcommand{\mZ}{\bm{Z}}

% Bold Greek matrices retained from the graph papers.
\newcommand{\mBeta}{\bm{\beta}}
\newcommand{\mDelta}{\bm{\Delta}}
\newcommand{\mGamma}{\bm{\Gamma}}
\newcommand{\mLambda}{\bm{\Lambda}}
\newcommand{\mOmega}{\bm{\Omega}}
\newcommand{\mPhi}{\bm{\Phi}}
\newcommand{\mPi}{\bm{\Pi}}
\newcommand{\mSigma}{\bm{\Sigma}}
\newcommand{\mTheta}{\bm{\Theta}}

% Calligraphic symbols.
\newcommand{\gA}{\mathcal{A}}
\newcommand{\gB}{\mathcal{B}}
\newcommand{\gC}{\mathcal{C}}
\newcommand{\gD}{\mathcal{D}}
\newcommand{\gE}{\mathcal{E}}
\newcommand{\gF}{\mathcal{F}}
\newcommand{\gG}{\mathcal{G}}
\newcommand{\gH}{\mathcal{H}}
\newcommand{\gI}{\mathcal{I}}
\newcommand{\gJ}{\mathcal{J}}
\newcommand{\gK}{\mathcal{K}}
\newcommand{\gL}{\mathcal{L}}
\newcommand{\gM}{\mathcal{M}}
\newcommand{\gN}{\mathcal{N}}
\newcommand{\gO}{\mathcal{O}}
\newcommand{\gP}{\mathcal{P}}
\newcommand{\gQ}{\mathcal{Q}}
\newcommand{\gR}{\mathcal{R}}
\newcommand{\gS}{\mathcal{S}}
\newcommand{\gT}{\mathcal{T}}
\newcommand{\gU}{\mathcal{U}}
\newcommand{\gV}{\mathcal{V}}
\newcommand{\gW}{\mathcal{W}}
\newcommand{\gX}{\mathcal{X}}
\newcommand{\gY}{\mathcal{Y}}
\newcommand{\gZ}{\mathcal{Z}}

% Blackboard-bold symbols.
\newcommand{\sA}{\mathbb{A}}
\newcommand{\sB}{\mathbb{B}}
\newcommand{\sC}{\mathbb{C}}
\newcommand{\sD}{\mathbb{D}}
\newcommand{\sE}{\mathbb{E}}
\newcommand{\sF}{\mathbb{F}}
\newcommand{\sG}{\mathbb{G}}
\newcommand{\sH}{\mathbb{H}}
\newcommand{\sI}{\mathbb{I}}
\newcommand{\sJ}{\mathbb{J}}
\newcommand{\sK}{\mathbb{K}}
\newcommand{\sL}{\mathbb{L}}
\newcommand{\sM}{\mathbb{M}}
\newcommand{\sN}{\mathbb{N}}
\newcommand{\sO}{\mathbb{O}}
\newcommand{\sP}{\mathbb{P}}
\newcommand{\sQ}{\mathbb{Q}}
\newcommand{\sR}{\mathbb{R}}
\newcommand{\sS}{\mathbb{S}}
\newcommand{\sT}{\mathbb{T}}
\newcommand{\sU}{\mathbb{U}}
\newcommand{\sV}{\mathbb{V}}
\newcommand{\sW}{\mathbb{W}}
\newcommand{\sX}{\mathbb{X}}
\newcommand{\sY}{\mathbb{Y}}
\newcommand{\sZ}{\mathbb{Z}}

% Frequently used operators.
\DeclareMathOperator{\Cov}{Cov}
\DeclareMathOperator{\Var}{Var}
\DeclareMathOperator{\diag}{diag}
\DeclareMathOperator{\nei}{Nei}
\DeclareMathOperator{\nnz}{nnz}
\DeclareMathOperator{\pr}{pr}
\DeclareMathOperator{\prox}{prox}
\DeclareMathOperator{\push}{push}
\DeclareMathOperator{\res}{res}
\DeclareMathOperator{\rel}{Rel}
\DeclareMathOperator{\relu}{ReLU}
\DeclareMathOperator{\sign}{sign}
\DeclareMathOperator{\supp}{supp}
\DeclareMathOperator{\tr}{tr}
\DeclareMathOperator{\Tr}{Tr}
\DeclareMathOperator{\vol}{vol}
\DeclareMathOperator{\work}{work}
\DeclareMathOperator{\Work}{Work}
\DeclareMathOperator*{\argmax}{arg\,max}
\DeclareMathOperator*{\argmin}{arg\,min}

% Project-wide names and claim-status labels used by standalone research notes.
% Mathematical notation belongs in math_commands.tex.

% Algorithms and solver families.
\newcommand{\AESP}{\textsc{AESP}}
\newcommand{\APPR}{\textsc{APPR}}
\newcommand{\ASPR}{\textsc{ASPR}}
\newcommand{\FISTA}{\textsc{FISTA}}
\newcommand{\ISTA}{\textsc{ISTA}}
\newcommand{\LOCAPPR}{\textsc{LocAPPR}}
\newcommand{\LOCGD}{\textsc{LocGD}}
\newcommand{\LOCSOR}{\textsc{LocSOR}}
\newcommand{\LOCCH}{\textsc{LocCH}}
\newcommand{\LOCHB}{\textsc{LocHB}}
\newcommand{\LocProxGS}{\textsc{LocProxGS}}
\newcommand{\Hybrid}{\textsc{AESP--LocSOR}}

% Claim classifications required by the repository research-note protocol.
\newcommand{\statussource}{\textbf{Source result}}
\newcommand{\statusproved}{\textbf{Proved here}}
\newcommand{\statusconditional}{\textbf{Conditional}}
\newcommand{\statusconjecture}{\textbf{Open / conjectural}}
\newcommand{\statusopen}{\textbf{Open}}
\newcommand{\statusempirical}{\textbf{Empirical}}
\newcommand{\statusobservation}{\textbf{Empirical observation}}
\newcommand{\statusrefuted}{\textbf{Refuted route}}

% Compatibility names used by the older complete-note presentation layer.
\newcommand{\Status}[2]{\textbf{#1:} #2}
\newcommand{\SourceStatus}{\textnormal{\textsc{Source result}}}
\newcommand{\ProvedStatus}{\textnormal{\textsc{Proved here}}}
\newcommand{\ConditionalStatus}{\textnormal{\textsc{Conditional}}}
\newcommand{\ComputedStatus}{\textnormal{\textsc{Computed}}}
\newcommand{\RefutedStatus}{\textnormal{\textsc{Refuted route}}}
\newcommand{\OpenStatus}{\textnormal{\textsc{Open}}}



\title{An Analytic Route to the Clipped Acceleration Master Inequality}
\author{Baojian Zhou\\
\small Working research note for \texttt{hybrid-local-solver}}
\date{August 2026}

\begin{document}
\maketitle

\begin{abstract}
This note independently reconstructs the exact one-step Lyapunov master form
for a clipped accelerated proximal recurrence.  The identity is a statement
about an exact shifted solve on one positive face; it is not a convergence or
local-work theorem.  The reconstruction exposes two sign-visible channels,
leaving only mixed clipping as the analytic obstruction.  Fable's finite
exact decisions are retained as source evidence rather than promoted here.
\end{abstract}

% Canonical source-aligned PageRank/RPPR reference formulation.
%
% This fragment is intentionally heading-free at the section level so that the
% active paper and every standalone note can input it. It is a manuscript
% reference convention, not an implementation-wide residual decision.
% Primary source: Fountoulakis and Martinez-Rubio, arXiv:2602.21138v2,
% PDF p. 3, Section 3 and equation (RPPR).

\subsection{Common source-aligned PageRank and RPPR model}
\label{subsec:shared-source-aligned-problem}

All documents in this manuscript workspace use the following reference
language. It follows the plain italic notation of the comparison paper:
\(x,s,A,D,Q,I\) denote column vectors or matrices as determined by their
displayed domains. This is the scoped exception to the project's usual
bold-vector and bold-matrix typography. A note may impose a stronger
assumption after this block, but it must not redefine these objects.

Let \(G=(V,E)\) be a finite, simple, undirected, unweighted graph without
isolated vertices, let \(V=[n]:=\{1,\ldots,n\}\), and let
\(A\in\{0,1\}^{n\times n}\) be its symmetric adjacency matrix. For
\(i\in V\), write
\[
    \mathcal N(i):=\{j\in V:\{i,j\}\in E\},
    \qquad
    d_i:=|\mathcal N(i)|,
    \qquad
    D:=\diag(d_1,\ldots,d_n).
\]
For \(S\subseteq V\), define
\[
    \mathcal N(S):=\bigcup_{i\in S}\mathcal N(i),
    \qquad
    \partial S:=\mathcal N(S)\setminus S,
    \qquad
    \operatorname{Ext}(S):=V\setminus(S\cup\partial S),
\]
\begin{equation}
    \vol(S):=\sum_{i\in S}d_i.
    \label{eq:shared-volume}
\end{equation}
The support of a vector is
\(\supp(x):=\{i\in[n]:x_i\neq0\}\). Matrix inequalities use the Loewner
order, and all unqualified vector inequalities are coordinatewise.

The seed is a column distribution
\begin{equation}
    s\in\R_+^n,
    \qquad
    \inner{\one}{s}=1.
    \label{eq:shared-seed}
\end{equation}
The single-seed specialization is always written \(s=e_v\), where \(v\in V\)
is the seed vertex. Thus \(s\) is never used as a vertex label. For
\(\alpha\in(0,1]\), define
\begin{equation}
\begin{aligned}
    \mathcal L&:=I-D^{-1/2}AD^{-1/2},\\
    Q&:=\alpha I+\frac{1-\alpha}{2}\mathcal L
      =\frac{1+\alpha}{2}I
       -\frac{1-\alpha}{2}D^{-1/2}AD^{-1/2},\\
    b&:=\alpha D^{-1/2}s.
\end{aligned}
    \label{eq:shared-pagerank-matrices}
\end{equation}
Since \(0\preceq\mathcal L\preceq2I\),
\(\alpha I\preceq Q\preceq I\). The shared smooth PageRank quadratic is
\begin{equation}
    f(x):=\frac12\inner{x}{Qx}-\inner{b}{x},
    \qquad
    \nabla f(x)=Qx-b.
    \label{eq:shared-pagerank-objective}
\end{equation}
Its unique unregularized minimizer and associated lazy personalized PageRank
vector are
\begin{equation}
    x^0:=Q^{-1}b,
    \qquad
    \pi:=D^{1/2}x^0.
    \label{eq:shared-ppr-solution}
\end{equation}

For a regularization parameter \(\rho>0\), reserve \(g_\rho\) for the
nonsmooth weighted \(\ell_1\) term and define
\begin{equation}
    g_\rho(x):=\alpha\rho\norm{D^{1/2}x}_1,
    \qquad
    F_\rho(x):=f(x)+g_\rho(x),
    \qquad
    x^\star(\rho):=\argmin_{x\in\R^n}F_\rho(x).
    \label{eq:shared-rppr-objective}
\end{equation}
When \(\rho\) is fixed, \(x^\star\) abbreviates \(x^\star(\rho)\), never the
unregularized point \(x^0\). Write
\[
    S^\star(\rho):=\supp(x^\star(\rho)),
    \qquad
    I^\star(\rho):=[n]\setminus S^\star(\rho).
\]
The source records \(x^\star(\rho)\geq0\),
\(\vol(S^\star(\rho))\leq1/\rho\), and the coordinatewise conditions
\begin{equation}
    \nabla_i f(x^\star(\rho))
    \in
    \begin{cases}
        \{-\alpha\rho\sqrt{d_i}\}, & x_i^\star(\rho)>0,\\
        [-\alpha\rho\sqrt{d_i},0], & x_i^\star(\rho)=0.
    \end{cases}
    \label{eq:shared-rppr-kkt}
\end{equation}

The common computational unit is one repeated adjacency-list scan: scanning
coordinate \(i\) costs \(d_i\), and scanning a set \(S\) costs \(\vol(S)\).
Iteration count and wall-clock time do not replace this degree-weighted work
measure.

Finally, accuracy symbols remain distinct. The APPR threshold is
\(\epsappr\); \(\epsobj\) denotes an objective-gap target;
\(\epspg\) denotes the source experiment's proximal fixed-point tolerance; and
\(\epsppr\) may denote a document-scoped degree-normalized PPR error target.
No equality or conversion among these quantities is assumed. In particular,
this shared model does not resolve the repository-wide residual decision.


\section{Scope, coordinates, and evidence labels}

The shared RPPR objective uses the normalized variable in the imported
definition above.  This note temporarily uses degree-scaled coordinates
$u=D^{-1/2}x$, because clipping by a common degree-normalized threshold then
acts as $(u-c\one)_+$.  In these coordinates the quadratic matrix is
\begin{equation}
 \widetilde Q
 :=D^{1/2}QD^{1/2}
 =\frac{1+\alpha}{2}D-\frac{1-\alpha}{2}A.
 \label{eq:psi-scaled-q}
\end{equation}
This is a note-scoped change of coordinates, not a change to the shared RPPR
convention.

\paragraph{Source.}
Fable's iteration-7 files propose the master form and report exact finite
decisions.  They are exploratory evidence, not proof authority for this note.
\paragraph{Proved here.}
The recurrence, admissible functional domain, master identity, and the two
unmixed-sign inequalities below are derived independently.
\paragraph{Open.}
Nonpositivity for mixed clipping on general graphs, proper-face transport,
finite shifted solves, and every end-to-end work conclusion remain open at
this point in the note.

\section{Exact recurrence and admissible functional domain}

Fix $0<q\leq1/2$ and put
\begin{equation}
 \alpha=\frac{q^2}{1+q^2},\qquad
 \kappa=1-2\alpha,\qquad
 \beta=\frac{1-q}{1+q},\qquad
 \nu=\frac{q}{1+q}=\frac{1-\beta}{2}.
 \label{eq:psi-parameters}
\end{equation}
Let $\langle u,v\rangle_D=u^TDv$, let
$\mathsf v_G=\one^TD\one$, and let
\begin{equation}
 P=I-\frac{\one\one^TD}{\mathsf v_G}
 \label{eq:psi-projection}
\end{equation}
be the $D$-orthogonal projection onto $\one^{\perp_D}$.  Define the exact
shifted-resolvent operator
\begin{equation}
 M:=\kappa(\widetilde Q+\kappa D)^{-1}D,
 \qquad
 m_0:=\frac{\kappa}{\kappa+\alpha}=1-q^2.
 \label{eq:psi-resolvent}
\end{equation}
The matrix $M$ is $D$-self-adjoint, positive definite, entrywise
nonnegative, commutes with $P$, and satisfies $M\one=m_0\one$.  The last
display and \cref{eq:psi-parameters} give the two scalar identities
\begin{equation}
 \beta m_0=(1-q)^2,
 \qquad
 m_0(1+\beta)^2=4\beta.
 \label{eq:psi-scalar-identities}
\end{equation}

Consider one exact positive-face shifted step.  With
$e_t=x^\star-x_t$, $d_t=x_t-x_{t-1}=e_{t-1}-e_t\geq0$, and an arbitrary
cap $\Delta_t\geq0$, put
\begin{equation}
 y_t=\beta d_t-\Delta_t\one,
 \qquad
 \phi_t=(y_t)_+,
 \qquad
 \ell_t=x_t+\phi_t.
 \label{eq:psi-clip}
\end{equation}
Exact minimization of the shifted quadratic centered at $\ell_t$ is
equivalent to
\begin{equation}
 e_{t+1}=M(e_t-\phi_t).
 \label{eq:psi-error-recurrence}
\end{equation}
This reconstruction assumes that the obstacle solve stays on the stated
positive face.  It does not silently extend to a proper face.

Write
\begin{equation}
 h_{t-1}=Pe_{t-1},\quad h_t=Pe_t,\quad
 z_t=Py_t=\beta(h_{t-1}-h_t),\quad
 w_t=P(y_t)_-,\quad b_t=z_t-\nu h_t.
 \label{eq:psi-high-variables}
\end{equation}
Since $(y_t)_+=y_t+(y_t)_-$, one has
$P\phi_t=z_t+w_t$ and therefore
\begin{equation}
 h_{t+1}=M(h_t-z_t-w_t).
 \label{eq:psi-high-recurrence}
\end{equation}

The trajectory-free functional below is defined for every
$y\in\mathbb R^V$ and $h\in\one^{\perp_D}$.  This really is the full
algebraic cap domain: given any $y$, choosing
$\Delta\geq\max_i(-y_i)$ and $d=(y+\Delta\one)/\beta$ gives
$d\geq0$ and $y=\beta d-\Delta\one$.  The functional domain is nevertheless
stronger than reachability by a single fixed RPPR trajectory, because it
forgets the preceding history and the positive-face inequalities.

\section{Independent master-identity reconstruction}

For $u,v\in\one^{\perp_D}$ define
\begin{equation}
 \mathcal V(u,v)
 :=\norm{u}_D^2-(1+\beta)\langle u,Mv\rangle_D
     +\beta\langle v,Mv\rangle_D.
 \label{eq:psi-lyapunov}
\end{equation}
Also put
\begin{equation}
 G:=m_0I-2M+M^2.
 \label{eq:psi-g-operator}
\end{equation}

\begin{lemma}[Exact clipped master identity]
\label{lem:psi-master-identity}
For every step satisfying \cref{eq:psi-error-recurrence},
\begin{equation}
 \mathcal V(h_{t+1},h_t)-(1-q)^2\mathcal V(h_t,h_{t-1})
 =\Psi(y_t,h_t),
 \label{eq:psi-master-identity}
\end{equation}
where, for arbitrary $y$ and $h\in\one^{\perp_D}$,
\begin{align}
 z&:=Py,& w&:=P(y_-),& b&:=z-\nu h,\notag\\
 \Psi(y,h)
 &:=\norm{M(b+w)}_D^2-m_0\langle b,Mb\rangle_D
       -\nu^2\norm{Mh}_D^2-\beta\langle h,Gh\rangle_D.
 \label{eq:psi-master-form}
\end{align}
\end{lemma}

\begin{proof}
Put $\bar\phi=P\phi=z+w$ and $c=b/\beta$.  From
$z=\beta(h_{t-1}-h_t)$ one has
$h_{t-1}=h_t+z/\beta$ and $z=\beta c+\nu h_t$.
Substitute these two relations and
$h_{t+1}=M(h_t-\bar\phi)$ into \cref{eq:psi-lyapunov}.  Using
$D$-self-adjointness of $M$ and the two scalar identities in
\cref{eq:psi-scalar-identities}, direct collection gives
\begin{align}
 &\mathcal V(h_{t+1},h_t)-(1-q)^2\mathcal V(h_t,h_{t-1})\notag\\
 &\quad=\norm{M\bar\phi}_D^2
 -(1-\beta)\langle Mh_t,M\bar\phi\rangle_D
 -\beta^2m_0\langle c,Mc\rangle_D
 -\beta\langle h_t,Gh_t\rangle_D.
 \label{eq:psi-collected-form}
\end{align}
Finally, $b=\beta c=z-\nu h_t$ and
$M(b+w)=M(\bar\phi-\nu h_t)$.  Expanding its squared norm and using
$1-\beta=2\nu$ turns \cref{eq:psi-collected-form} exactly into
\cref{eq:psi-master-form}.
\end{proof}

The calculation uses neither the trigger that selected $\Delta_t$ nor a
stage-type case split.  Conversely, it uses exact shifted minimization and a
fixed positive face; it is not an identity for a finite inner solve without
additional residual terms.

\section{A sign-visible decomposition}

Expanding only the first square in \cref{eq:psi-master-form} gives
\begin{align}
 \Psi(y,h)
 &=\norm{Mw}_D^2+2\langle M^2w,b\rangle_D
   -\langle b,M(m_0I-M)b\rangle_D\notag\\
 &\qquad-\nu^2\norm{Mh}_D^2-\beta\langle h,Gh\rangle_D.
 \label{eq:psi-slack-decomposition}
\end{align}
This representation keeps the complementarity $y_+y_-=0$ visible; replacing
$y_+$ and $y_-$ by independent vectors would discard the main structure.

Let $\mu_2$ be the first nonzero normalized-Laplacian eigenvalue.  On a mode
with normalized-Laplacian eigenvalue $\mu$, the generalized eigenvalue of
$\widetilde Q$ is
$\lambda(\mu)=\alpha+(1-\alpha)\mu/2$, and the eigenvalue of $M$ is
$m(\mu)=\kappa/(\kappa+\lambda(\mu))$.  If
\begin{equation}
 \mu_2\geq2q,
 \label{eq:psi-hgap}
\end{equation}
then $0<m(\mu)\leq1-q<m_0$ on $\one^{\perp_D}$.  Since
$m_0-2m+m^2\geq0$ for $0<m\leq1-q$, both
$M(m_0I-M)$ and $G$ are positive semidefinite on the high space.

Two entire channels now have a visible sign:
\begin{itemize}
 \item If $y\geq0$, then $w=0$, and every term in
 \cref{eq:psi-slack-decomposition} is nonpositive.
 \item If $y\leq0$, then $w=-z$ and $b+w=-\nu h$, so
 \begin{equation}
  \Psi(y,h)=-m_0\langle b,Mb\rangle_D
            -\beta\langle h,Gh\rangle_D\leq0.
 \end{equation}
\end{itemize}
Thus only a genuinely mixed sign pattern can violate the master inequality.
This conclusion is analytic and family-independent under
\cref{eq:psi-hgap}; it does not rely on finite enumeration.

\section{A structural mixed-clipping theorem}

The remaining mixed term admits a sufficient condition with a probabilistic
sign interpretation.  Assume
\begin{equation}
 (M^2)_{ij}\leq \frac{m_0^2d_j}{\mathsf v_G}
 \qquad(i\neq j).
 \label{eq:psi-hk}
\end{equation}
This condition is entrywise and stronger than the desired conclusion, but it
can be checked analytically on graph families.

\begin{lemma}[Kernel positive association]
\label{lem:psi-kernel-positive-association}
Under \cref{eq:psi-hk}, for every $a\in\mathbb R^V$ and threshold
$\Delta\in\mathbb R$, if
\begin{equation*}
 f=(a-\Delta\one)_+,
 \qquad
 g=a-f=\min\{a,\Delta\one\},
\end{equation*}
then
\begin{equation}
 \langle MPf,MPg\rangle_D\geq0.
 \label{eq:psi-positive-association}
\end{equation}
\end{lemma}

\begin{proof}
Let $\pi_i=d_i/\mathsf v_G$ and define
\begin{equation*}
 K_{ij}:=\frac{d_i(M^2)_{ij}}{m_0^2\mathsf v_G}.
\end{equation*}
Entrywise nonnegativity of $M$, $D$-self-adjointness, and
$M^2\one=m_0^2\one$ show that $K$ is a symmetric probability coupling with
both marginals $\pi$.  Condition \cref{eq:psi-hk} is precisely
$K_{ij}\leq\pi_i\pi_j$ off the diagonal.  Centering by $P$ and using the
marginals gives the exact two-point identity
\begin{align}
 \langle MPf,MPg\rangle_D
 &=m_0^2\mathsf v_G
   \left(\sum_{i,j}K_{ij}f_ig_j
        -\sum_i\pi_if_i\sum_j\pi_jg_j\right)\notag\\
 &=\frac{m_0^2\mathsf v_G}{2}
   \sum_{i,j}(\pi_i\pi_j-K_{ij})(f_i-f_j)(g_i-g_j).
 \label{eq:psi-two-point-kernel}
\end{align}
Both $f$ and $g$ are nondecreasing functions of the same scalar $a_i$, so
the difference product is nonnegative.  Its diagonal value is zero, and
every off-diagonal coefficient is nonnegative by \cref{eq:psi-hk}.  This
proves the claim.
\end{proof}

\begin{theorem}[Structural nonpositivity of the master form]
\label{thm:psi-hk-nonpositivity}
If $\mu_2\geq2q$ and \cref{eq:psi-hk} holds, then
\begin{equation}
 \Psi(y,h)\leq0
 \qquad
 \text{for every }y\in\mathbb R^V,
 \quad h\in\one^{\perp_D}.
 \label{eq:psi-hk-conclusion}
\end{equation}
\end{theorem}

\begin{proof}
Given $y$, choose $\Delta\geq\max_i(-y_i)$ and put
$a=y+\Delta\one$.  Then $f=y_+$, $z=Py=Pa$, and
$Pg=z-Pf$.  Hence \cref{lem:psi-kernel-positive-association} says
\begin{equation}
 \langle MPf,M(z-Pf)\rangle_D\geq0,
 \quad\text{or equivalently}\quad
 \langle MPf,Mz\rangle_D\geq\norm{MPf}_D^2.
 \label{eq:psi-kernel-payment}
\end{equation}
In \cref{eq:psi-collected-form}, write $c=b/\beta$ and note
$z=\beta c+\nu h$.  The first two terms of that form satisfy
\begin{align*}
 L
 &:=\norm{MPf}_D^2-2\nu\langle Mh,MPf\rangle_D\\
 &=\norm{MPf}_D^2-2\langle M(z-\beta c),MPf\rangle_D\\
 &\leq-\norm{MPf}_D^2+2\beta\langle Mc,MPf\rangle_D\\
 &\leq\beta^2\norm{Mc}_D^2
 \leq\beta^2m_0\langle c,Mc\rangle_D.
\end{align*}
The first inequality uses \cref{eq:psi-kernel-payment}, the second completes
the square, and the third uses $0\prec M\preceq m_0I$ in the
$D$-inner product.  This pays the $c$-term of
\cref{eq:psi-collected-form}.  The remaining term
$-\beta\langle h,Gh\rangle_D$ is nonpositive under
$\mu_2\geq2q$, proving \cref{eq:psi-hk-conclusion}.
\end{proof}

\section{An infinite non-scalar family: complete bipartite graphs}

Condition \cref{eq:psi-hk} has a $q$-independent graph form.  Indeed,
\begin{equation}
 \widetilde Q+\kappa D
 =\frac{1-\alpha}{2}(3D-A),
 \qquad
 M=\frac{2(1-2\alpha)}{1-\alpha}N,
 \qquad
 N:=(3D-A)^{-1}D.
 \label{eq:psi-n-resolvent}
\end{equation}
Since $N\one=\one/2$, condition \cref{eq:psi-hk} is equivalent to
\begin{equation}
 4\mathsf v_G(N^2)_{ij}\leq d_j
 \qquad(i\neq j).
 \label{eq:psi-hk-n}
\end{equation}

\begin{theorem}[Master nonpositivity on every complete bipartite graph]
\label{thm:psi-complete-bipartite}
Let $G=K_{a,b}$ with integers $a,b\geq1$.  For every $0<q\leq1/2$,
condition \cref{eq:psi-hgap} holds and
\begin{equation}
 \sup_{y\in\mathbb R^{a+b},\ h\in\one^{\perp_D}}\Psi(y,h)=0.
 \label{eq:psi-complete-bipartite-sup}
\end{equation}
In particular, the exact positive-face recurrence satisfies the one-step
contraction \cref{eq:psi-master-identity} on the entire infinite family.
\end{theorem}

\begin{proof}
Let $L$ and $R$ be the two parts, with $|L|=a$ and $|R|=b$; degrees are $b$
on $L$ and $a$ on $R$, and $\mathsf v_G=2ab$.  Solving
$(3D-A)z=Dx$ by the two part sums gives
\begin{equation}
 N=
 \begin{pmatrix}
 \frac13I_a+\frac1{24a}J_a&\frac1{8b}J_{a\times b}\\[2mm]
 \frac1{8a}J_{b\times a}&\frac13I_b+\frac1{24b}J_b
 \end{pmatrix}.
 \label{eq:psi-kab-n}
\end{equation}
Squaring the four blocks shows, for distinct vertices,
\begin{equation}
 (N^2)_{ij}=
 \begin{cases}
  13/(288a),&i,j\in L,\\
  13/(288b),&i,j\in R,\\
  3/(32b),&i\in L,\ j\in R,\\
  3/(32a),&i\in R,\ j\in L.
 \end{cases}
 \label{eq:psi-kab-n2}
\end{equation}
The corresponding left-to-right ratios in \cref{eq:psi-hk-n} are
$13/36$ within either part and $3/4$ across the cut.  Thus
\cref{eq:psi-hk-n} holds strictly.

The normalized-Laplacian spectrum of $K_{a,b}$ is
$\{0,1^{(a+b-2)},2\}$ when $a+b\geq3$, while $K_{1,1}$ has spectrum
$\{0,2\}$.  Hence $\mu_2\geq1\geq2q$.  The structural theorem
\cref{thm:psi-hk-nonpositivity} gives $\Psi\leq0$.  Equality is attained at
$y=0,h=0$, so the supremum is zero.
\end{proof}

This family is not a disguised scalar-high-space case: when $a+b\geq3$, the
normalized Laplacian has the two distinct high eigenvalues $1$ and $2$.
Thus the theorem goes beyond the complete-graph truncation-covariance proof.
It includes every star $K_{1,b}$, but it remains a positive-face Lyapunov
statement rather than a locality or work theorem.

\section{Current boundary}

The positive-association route is now closed on every complete bipartite
graph.  The next falsifiable targets are: extend the block-resolvent
calculation to complete multipartite or other equitable-partition families;
or find an exact graph satisfying $\mu_2\geq2q$ but violating
\cref{eq:psi-hk} and then decide $\Psi$ analytically there.  A one-step
inequality implies Lyapunov absorption only for the exact positive-face
recurrence above.  It does not by itself charge graph discovery, clipping,
shifted inverse work, proper-face changes, finite precision, or final PPR
output.

\end{document}
